\documentclass{visionen}
\ArticleStart{recto}

% enable bold text globally:
\let\bfseries\DesignBoldSeries


\begin{document}

% this is the default title format for a two-column article

\ArticleHeading{LaTeX Article Template}{Your Name - and whatever you want to add to your byline}

% uncomment the code below to enable a whole-page-spanning title format instead, similar to the "classic" VISIONEN article layout.

%\twocolumn[
% \ArticleHeadingWide{LaTeX Article Template for the Longer Titles Among Them All}{Your Name - and whatever you want to add to your byline}
%]

% if you want to reduce space here, use \vspace{-XXpt} (replace XX with 12-30)

\OpeningParagraph{Start with a title and byline, then put the opening paragraph inside this command. It supplies the three-line initial and tracked opening words. The source of this short article demonstrates every shared feature.}

Paragraphs are generally indented. However, the opening paragraph, as well the first paragraph after each heading (like the paragraph immediately below the next heading) is not.

% if you don't want your sections to be numbered, use \section*{Section Name}
\section{Limitations}

Try to use \emph{italics} for text that you want to emphasize. Depending on the context, you could also use \textsc{small capitals}. \textbf{Bold} text for this font is rather thick, so don't use it unless you want to highlight something massive, not more than once per text-rich page (best: not more than once per spread).

The title of the article is limited by default to only the left column of a two-column layout. In case you haven't tried it out yet, try commenting out the current article heading code and uncommenting the wide article heading code that comes directly after. 

\subsection{How to work with this template}

\begin{enumerate}
  \item The class should be kept in the same folder as the article text. Fonts and images should be placed in a separate folder (see template folder structure).
  \item Run \texttt{lualatex example.tex}. Alternatively, you can compile the document on Overleaf, just make sure to change the compiler to LuaLaTex under \texttt{File > Settings > Compiler}.
\end{enumerate}

\begin{itemize}
  \item Lists have compact spacing.
  \item The image wrapper reserves whole grid rows.
\end{itemize}

\noindent Shared TikZ styles supply vertex type, rounded strokes, and edge clearances:

\GridBlock{\begin{tikzpicture}[math graph,x=24bp,y=18bp]
  \node (a) at (0,0) {$a$};
  \node (b) at (2,0) {$b$};
  \node (c) at (1,1) {$c$};
  \draw[graph edge] (a) -- (c) -- (b);
  \draw[graph full edge] (a) -- (b);
  \draw[relation rule] (-0.4,-0.5) -- (2.4,-0.5);
\end{tikzpicture}}

% i don't think anyone will use this but ok 
\begin{Caption}{9bp}
A three-vertex graph, with a 9 bp inset.\\*
The image caption below uses no inset.
\end{Caption}
\pagebreak[0]

% add an image here to see how it will look
\GridBlock{\includegraphics[width=0.9\columnwidth]{images/put_images_into_images_folder.png}}

\begin{Caption}{0bp}
You can add captions to your images. Captions use an inset of 0 bp; verse can use 9 bp.
\end{Caption}

\section{Math and code}

Write $f(x) = x^2$ inline. To display math in the center of the text body, you can follow this exmple:

\GridMath{
  \begin{aligned}
    S_n &= 1 + 2 + \cdots + n \\
        &= \frac{n(n+1)}{2}.
  \end{aligned}
}

\noindent The CodeBlock class can be used to write code. The class supplies the appropriate font, theme, and wrapping:

\begin{CodeBlock}{python}
def triangular(n: int) -> int:
    return n * (n + 1) // 2
\end{CodeBlock}

\noindent Choose another language:

\begin{CodeBlock}{javascript}
const triangular = n => n * (n + 1) / 2;
\end{CodeBlock}

\section{Page-level layout}

You might have noticed that this page has a larger left margin than the previous one. It's because it's intended to be the right-side page. Generally, each article gets compiled with the assumption that it starts on the left-side page of a spread. However, as these LaTeX pages ultimately get pasted into Affinity, this margin can be re-adjusted after your article is finished anyway. So don't worry about this.

\section*{Finishing with sources}

Each source should be listed separately. The DOI command adds a link.

\begin{GridReferences}
  \item Lipski and Preparata (1981).
  \textit{Efficient Algorithms for Finding Maximum Matchings in Convex Bipartite Graphs and Related Problems}.
  \DOI{10.1007/bf00264533}
\end{GridReferences}

\end{document}
